\usepackage{verbatim}
\usepackage{listings} % Para codigo fuente y trozos de codigo
\lstloadlanguages{C,bash}
\lstset{language=C, 
        basicstyle=\scriptsize\ttfamily,
        breaklines=true,
        keywordstyle=\color{blue}\ttfamily,
        stringstyle=\color{red}\ttfamily,
        commentstyle=\color{gray}\ttfamily,
        morecomment=[l][\color{magenta}]{\#},
        numbers=left, numberstyle=\tiny
}
\usepackage{minted}
\usepackage[spanish,es-tabla]{babel}
\usepackage{siunitx}
\usepackage{graphicx}
\graphicspath{{capitulos/figuras/}}
\usepackage[]{caption} %labelsep=period
\usepackage{float}

\usepackage[
backend=biber,
style=apa,
]{biblatex}
%%% Editar %%%
\addbibresource{references.bib}
\let\cite\parencite
%%%%%%%%%%%%%%

\usepackage{fontspec}
\input{esConfigDir/esConfigFonts}
\usepackage{authblk}
\DefineBibliographyStrings{spanish}{andothers={\itshape{et~al}\adddot}}

%In a current version of biblatex-apa, the comma before the "and" is either by the Oxford comma (\finalnamdelim)
%or produced by the internal macro \apablx@ifrevnameappcomma.
%If you want to get rid of the comma, you need to disable both of these features.
%https://tex.stackexchange.com/questions/630515/biblatex-remove-commas-after-author-first-names-in-references#:~:text=1%20Answer&text=In%20a%20current%20version%20of,disable%20both%20of%20these%20features.

\makeatletter
\renewcommand*{\apablx@ifrevnameappcomma}[2]{#2}
\makeatother
\DefineBibliographyExtras{spanish}{%
  \let\finalandcomma=\empty
}
% Fechas estilo APA en español: ordinal 1.º--9.º vía \mkbibmascord; día ≥ 10 sin ordinal.
% (No añadir punto antes de \sptext{o}: babel ya lo incluye en el ordinal.)
\DefineBibliographyExtras{spanish}{%
  \protected\def\bibapadayprint#1{%
    \ifnumgreater{\thefield{#1}}{9}
      {\stripzeros{\thefield{#1}}}
      {\mkbibmascord{\thefield{#1}}}%
  }%
  \protected\def\mkbibdateapalong#1#2#3#4{%
    \ifboolexpr{ (test {\ifentrytype{article}}
                  or (test {\ifentrytype{inbook}} and not test {\ifnameundef{editor}}))
                 and not test {\iffieldequalstr{entrysubtype}{nonacademic}} }
      {\clearfield{labelmonth}%
       \clearfield{labelday}}
      {}%
    \iffieldundef{#1}
      {}%
      {\iffieldbibstring{#1}{\bibncpstring{\thefield{#1}}}{\thefield{#1}}%
       \ifboolexpr{test {\iffieldundef{#4}} and test {\iffieldundef{#3}}}
         {}
         {\addcomma\space}}%
    \iffieldundef{#2}
      {}%
      {\iffieldundef{#1}
        {}%
        {\addcomma\addspace}%
       \iffieldbibstring{#2}{\bibcplstring{\thefield{#2}}}{\thefield{#2}}}%
    \iffieldundef{#4}
      {}%
      {\bibapadayprint{#4}%
       \iffieldundef{#3}
         {}
         {\addspace de\space}}%
    \iffieldundef{#3}
      {}%
      {\mkbibmonth{\thefield{#3}}}}%
  \protected\def\mkbibdateapalongextra#1#2#3#4{%
    \ifboolexpr{ (test {\ifentrytype{article}}
                  or (test {\ifentrytype{inbook}} and not test {\ifnameundef{editor}}))
                 and not test {\iffieldequalstr{entrysubtype}{nonacademic}} }
      {\clearfield{labelmonth}%
       \clearfield{labelday}}
      {}%
    \iffieldundef{#1}
      {}%
      {\iffieldbibstring{#1}{\bibncpstring{\thefield{#1}}}{\thefield{#1}}\printfield{extradate}%
       \ifboolexpr{test {\iffieldundef{#4}} and test {\iffieldundef{#3}}}
         {}
         {\addcomma\space}}%
    \iffieldundef{#2}
      {}%
      {\iffieldundef{#1}
        {}%
        {\addcomma\addspace}%
       \iffieldbibstring{#2}{\bibcplstring{\thefield{#2}}}{\thefield{#2}}}%
    \iffieldundef{#4}
      {}%
      {\bibapadayprint{#4}%
       \iffieldundef{#3}
         {}
         {\addspace de\space}}%
    \iffieldundef{#3}
      {}%
      {\mkbibmonth{\thefield{#3}}}}%
  \protected\def\mkbibdateapalongmdy#1#2#3#4{%
    \ifboolexpr{ (test {\ifentrytype{article}}
                  or (test {\ifentrytype{inbook}} and not test {\ifnameundef{editor}}))
                 and not test {\iffieldequalstr{entrysubtype}{nonacademic}} }
      {\clearfield{labelmonth}%
       \clearfield{labelday}}
      {}%
    \iffieldundef{#2}
      {}%
      {\iffieldundef{#1}
        {}%
        {\addcomma\addspace}%
       \iffieldbibstring{#2}{\bibcplstring{\thefield{#2}}}{\thefield{#2}}}%
    \iffieldundef{#4}
      {}%
      {el\addspace\bibapadayprint{#4}}%
    \iffieldundef{#3}
      {\addspace{en}\addspace}%
      {\iffieldundef{#4}
        {en\addspace}%
        {\addspace{de}\addspace}%
       \mkbibmonth{\thefield{#3}}{\addspace{de}\addspace}}%
    \iffieldundef{#1}
      {}%
      {\iffieldundef{#2}
        {}%
        {\addspace}%
       \iffieldbibstring{#1}{\bibstring{\thefield{#1}}}{\thefield{#1}}}}%
}
%Para cambiar el & por "y" en el capitulo de Referencias
\DeclareDelimFormat[bib,textcite]{finalnamedelim}{%
  \ifnumgreater{\value{liststop}}{2}{}{}%
  \addspace\bibstring{and}\space}

% Reemplaza las comillas angulares por comillas dobles normales
\usepackage{csquotes}
\DeclareQuoteStyle{spanish}{``}{''}{`}{'}

%para cambiar & en las citaciones dentro del texto
\DeclareDelimFormat[parencite]{finalnamedelim}{\addspace{y}\space}

\usepackage[letterspace=5]{microtype}
\usepackage{graphicx}

%%% Editar %%%
\def\texttitulocorto{Evaluación de una red IoT bajo el estándar LoRa}
%Lista de autores
\def\loauthors{J.\,N.~Poveda, L.\,M.~Santamaría, E.\,E.~Gaona García, S.\,Tabares, A.\,F.~Bustos}
%%%%%%%%%%%%%%
\input{esConfigDir/esEstilo}

\usepackage[
  paperheight=24.0cm,
  paperwidth=17.0cm,
  textwidth=12.5cm,
  tmargin=2.5cm,
  bmargin=2.5cm,
  left=2.5cm,
  right=2.0cm
]{geometry}


%%% Editar %%%
%Algunos paquetes para el usuario
\decimalpoint
\usepackage{lipsum}  
\usepackage[pangram]{blindtext}
%%%%%%%%%%%%%%


%%% Editar %%%
%% Comentar esta línea para la presentación si hipertexto
\usepackage[unicode,bookmarksopen,bookmarksdepth=2,breaklinks=true]{hyperref}%Para hipervinulos de figuras, tablas, ecuaciones y citas
\hypersetup{
        colorlinks=true,
        linkcolor=black,
        filecolor=black,
        urlcolor=black,
        citecolor=black
    }
% Tras hyperref: más puntos de ruptura en \url (p. ej. Dudler, R. en bibliografía).
% nobiblatex: evita que xurl deje las penalidades de biblatex demasiado bajas.
\usepackage[nobiblatex]{xurl}
\usepackage{etoolbox}
\gappto{\UrlBreaks}{\UrlOrds}
% URLs en bibliografía: permitir cortes en más posiciones.
\setcounter{biburlbreakpenalty}{100}
\setcounter{biburlnumpenalty}{9000}
\setcounter{biburlucpenalty}{9000}
\setcounter{biburllcpenalty}{9000}
% Campo url (biblatex-apa): no partir la secuencia tipográfica "https:"; el enlace sigue siendo la URL completa.
\ExplSyntaxOn
\str_new:N \l_bib_https_url_str
\cs_new_protected:Npn \bib_url_https_nobreak_print:n #1
  {
    \str_set:Nn \l_bib_https_url_str {#1}
    \str_if_eq:eeTF
      { \str_range:Nnn \l_bib_https_url_str { 1 } { 6 } }
      { https: }
      {
        \mbox{https:}
        \exp_args:Ne \nolinkurl
          {
            \str_range:Nnn \l_bib_https_url_str { 7 } { \str_count:N \l_bib_https_url_str }
          }
      }
      { \nolinkurl {#1} }
  }
\DeclareFieldFormat{url}{ \href{#1}{ \bib_url_https_nobreak_print:n {#1} } \nopunct }
\ExplSyntaxOff
%%%%%%%%%%%%%%

\usepackage{amsmath}
\usepackage{longtable}
\usepackage{subcaption}
\usepackage{array}
\newcolumntype{P}[1]{>{\centering\arraybackslash}p{#1}}

\usepackage{shapepar}
\usepackage{tikz}
\usepackage{tikzscale}
\usepackage{xcolor}
\definecolor{orcidlogocol}{HTML}{A6CE39}

\usepackage{tabularray}
\usepackage{nameref}
\usepackage{fontawesome}
\usepackage{academicons}

%% Para editar: comentar esta línea.
%% Para produccion del libro: activar esta linea.
%% Nota: Esta instrucción es para colocar lineas guia de corte (crop marks)
%\usepackage[cam,width=190truemm,height=260truemm,lualatex,center]{crop}
%%%%%%%%%%

%%% Editar %%%
%%%%%%%%%%%%%%

\usepackage[acronym,toc]{glossaries}
\usepackage{amsmath,amssymb,amsfonts,latexsym,cancel}
\usepackage{multirow}
\usepackage{booktabs}
\usepackage[table]{xcolor}
\usepackage{tikz} % hacer gráficos 
\usepackage{setspace}
\usepackage{xspace}

\def\fuente#1{\small\textbf{Fuente:}~#1}
\date{\today}

\newglossarystyle{long-tabla}{%
  \setglossarystyle{long}

  \renewcommand{\glossentry}[2]{%
    \small
    \glsentryitem{##1}\glstarget{##1}{\textmed{\glossentryname{##1}}} &
    \textls[-10]{\Glossentrydesc{##1}}\glspostdescription\space ##2\tabularnewline
  }%
}

\newglossarystyle{long-lista}{%
  \setglossarystyle{list}
  \renewcommand*{\glsgroupskip}{}%
  \renewcommand*{\glossentry}[2]{%
    \small
    \item[\glstarget{##1}{\textmed{\glossentryname{##1}}}] 
    \textls[-10]{\Glossentrydesc{##1}}\glspostdescription\space ##2
  }%
}

\setglossarystyle{long-lista}

%%% Sin sangria al inicio de párrafos
\setlength\parindent{0pt}
\raggedbottom
%%%

\pagestyle{empty}
\counterwithout{footnote}{chapter}

\makeglossaries

\def\texttitulo{Evaluación de una red IoT bajo el estándar LoRa}
\def\textsubtitulo{}
\def\loauthors{José~Noé~Poveda~Zafra, Luis~Enrique~Martín~Santamaría, Elvis~Eduardo~Gaona~García, Sebastián~Yesid~Tabares~Amaya, Alex~Francisco~Bustos~Pinzón}
\def\texttitulocorto{\texttitulo}

\usepackage{emptypage}

% Esto asegura que la página izquierda de un capítulo nuevo esté vacía (sin cornisas)
\renewcommand{\cleardoublepage}{%
  \clearpage\ifodd\value{page}\else\hbox{}\thispagestyle{empty}\newpage\fi%
}

\usepackage[all]{nowidow}

\newcommand{\CO}{CO\textsubscript{2}}
\newcommand{\IC}{I\textsuperscript{2}C}
\newcommand{\IS}{I\textsuperscript{2}S}
\newcommand{\NOX}{NO$_{x}$}
